\raggedbottom
Over the past decades, the auditing profession has undergone profound change.

When the first edition of \emph{Statistiek voor Accountants} appeared, statistical methods were already an important part of the auditor’s toolbox, but their practical application was often constrained by available technology, data accessibility, and computational effort. Calculating confidence intervals, determining sample sizes, evaluating regression models, or projecting sample results frequently required considerable manual work and specialized expertise. At the time, statistical sampling was experiencing renewed interest within the profession, particularly in areas such as compliance audits, securitisation reviews, and valuation engagements. The ability to understand statistical methods and critically evaluate their results was considered an essential competency for accountants and auditors.

\bigskip
Today, the environment in which auditors operate is fundamentally different.

Organizations generate and store vast amounts of data. Enterprise systems record transactions in real time. Analytical software has become widely available. Statistical analyses that once required substantial effort can now be performed in seconds. More recently, artificial intelligence has begun to automate not only computations, but also programming, visualization, and aspects of analytical interpretation.

Yet these developments have not diminished the importance of statistics in auditing. On the contrary, they have shifted the emphasis from performing analyses to understanding them.

As software increasingly performs the mechanics of statistical analysis, auditors must devote greater attention to questions such as:

\begin{itemize}
\item Is the selected method appropriate?
\item Which assumptions underlie the analysis?
\item How reliable are the results?
\item What uncertainty remains?
\item What conclusions are supported by the available evidence?
\item Which conclusions are not?
\end{itemize}

The responsibility for answering these questions remains firmly with the auditor.

\bigskip
This observation forms the foundation of the \emph{Audit Data Analysis} series.
The objective of these books is not simply to teach statistical techniques. Numerous software packages can perform statistical calculations, and an increasing number of AI tools can assist in their implementation. The enduring value of statistical education lies in developing the ability to understand, evaluate, and appropriately use statistical evidence.

\bigskip
For that reason, the series is built around three interconnected competencies:

\begin{itemize}
\item \textbf{Technical Skills} – the ability to perform statistical analyses correctly;
\item \textbf{Statistical Reasoning} – the ability to understand, interpret, and justify statistical methods;
\item \textbf{Professional Judgment} – the ability to evaluate evidence and determine which conclusions are supported by that evidence.
\end{itemize}

Together, these competencies support statistically informed professional judgment.

\bigskip
The series consists of two complementary volumes.

\textbf{Volume 1} concentrates on the statistical foundations of audit data analysis. Topics such as probability distributions, estimation, sampling, regression analysis, and goodness-of-fit techniques are introduced not as isolated mathematical subjects, but as tools for generating and evaluating evidence.

\textbf{Volume 2} focuses on the use of statistical evidence within the auditing process. Building on the foundations established in Volume 1, it addresses the application of analytical procedures and audit sampling within the framework of professional auditing standards.

The distinction between the two volumes reflects an important educational principle.
Before an auditor can determine how to respond to evidence, the auditor must first understand how evidence is generated, how its quality should be assessed, and what conclusions it supports.

In summary:

\begin{center}
\begin{tabular}{l l}
\textbf{Volume 1} & Understanding and evaluating statistical evidence \\
\textbf{Volume 2} & Using statistical evidence in auditing
\end{tabular}
\end{center}

Although the tools available to auditors have changed dramatically since the first editions of \emph{Statistiek voor Accountants}, the need for a sound understanding of statistical evidence remains unchanged. Indeed, developments in technology, data availability, and artificial intelligence have increased the importance of being able to critically evaluate analytical results rather than merely produce them. The ability to perform an analysis remains valuable, but the ability to understand its implications is increasingly essential.

The approach adopted in this series reflects this development. Statistical methods are presented not as ends in themselves, but as instruments for obtaining evidence. Throughout both volumes, emphasis is placed not only on performing analyses correctly, but also on understanding their assumptions, limitations, and implications.

We hope that this series helps students develop a deeper under\-standing of statistical evidence and its role within auditing. Equally, we hope it provides educators and practitioners with a framework that connects statistical methodology to the professional challenges faced by auditors in an increasingly data-rich environment.

\bigskip
August 2026,

Lucas Hoogduin

Paul Touw